\begin{theorem}[折叠算子的同余刚性与泛性质唯一性]\label{thm:fold-congruence-universal-property}
\leanverified{paper\_fold\_congruence\_universal\_property}
设 \(m\ge 1\)，\(\Omega_m:=\{0,1\}^m\)，并沿用定义 \ref{def:fold-word} 的加权值 \(N(\omega)=\sum_{k=1}^m \omega_kF_{k+1}\) 与折叠 \(\Fold_m:\Omega_m\to X_m\)。
定义值模映射
\[
r(\omega):=N(\omega)\bmod F_{m+2}\ \in\ \ZZ/(F_{m+2}\ZZ),
\]
并对唯一代表 \(0\le r\le F_{m+2}-1\) 定义
\[
w(r):=\pi_m\bigl(Z(r)\bigr)\in X_m.
\]
若映射 \(\Phi:\Omega_m\to X_m\) 满足对一切 \(\omega\in\Omega_m\)
\[
V_m(\Phi(\omega))\equiv N(\omega)\pmod{F_{m+2}},
\]
则必有 \(\Phi=\Fold_m\)。
此外，任意满足 \(\Phi|_{X_m}=\mathrm{id}_{X_m}\) 的回缩 \(\Phi:\Omega_m\to X_m\)，若同时满足上述同余条件，则亦唯一且等于 \(\Fold_m\)。
\end{theorem}

\begin{proof}
由引理 \ref{lem:zeckendorf-range-bijection}，映射 \(V_m:X_m\to\{0,1,\dots,F_{m+2}-1\}\) 为双射。
因此对任意 \(\omega\in\Omega_m\)，元素 \(V_m(\Phi(\omega))\) 落在 \(\{0,1,\dots,F_{m+2}-1\}\) 且与 \(N(\omega)\) 在模 \(F_{m+2}\) 意义下同余，从而必有
\[
V_m(\Phi(\omega))=r(\omega).
\]
又按 \(w(r)\) 的定义与引理 \ref{lem:zeckendorf-range-bijection}，\(V_m(w(r))=r\)，故由双射性得到
\(
\Phi(\omega)=w(r(\omega))
\)。
另一方面，由引理 \ref{lem:pom-fold-congruence}，\(\Fold_m(\omega)\) 仅依赖于 \(r(\omega)\)，并且 \(\Fold_m(w(r))=w(r)\)（命题 \ref{prop:fold-basic} 的幂等性）。
于是
\[
\Fold_m(\omega)=\Fold_m(w(r(\omega)))=w(r(\omega))=\Phi(\omega),
\]
从而 \(\Phi=\Fold_m\)。
最后，若 \(\Phi|_{X_m}=\mathrm{id}_{X_m}\)，则上证并未使用额外条件，仍推出 \(\Phi=\Fold_m\)，故唯一性成立。证毕。
\end{proof}

\begin{proposition}[Zeckendorf 补表示对偶：稳定类型上的内禀 \texorpdfstring{$\ZZ_2$}{Z2} 自同构与多重度守恒]\label{prop:fold-complement-duality}
\leanverified{paper\_fold\_complement\_duality}
定义补位对合 \(C:\Omega_m\to\Omega_m\) 为
\(
(C(\omega))_k:=1-\omega_k
\)。
并在模环 \(\ZZ/(F_{m+2}\ZZ)\) 上定义对合
\[
\tau_m(r)\equiv (F_{m+1}-2)-r\pmod{F_{m+2}}.
\]
由 \(w(r)=\pi_m(Z(r))\) 诱导 \(X_m\) 上的对合
\[
\jmath_m:X_m\to X_m,\qquad \jmath_m(w(r)):=w(\tau_m(r)).
\]
则：
\begin{enumerate}
\normalfont
  \item 对一切 \(\omega\in\Omega_m\) 有对偶相容性
  \[
  \Fold_m(C(\omega))=\jmath_m(\Fold_m(\omega)).
  \]
  \item 若定义纤维多重度 \(d_m(x):=\card{\Fold_m^{-1}(x)}\)，则对一切 \(x\in X_m\) 有守恒律
  \[
  d_m(x)=d_m(\jmath_m(x)).
  \]
\end{enumerate}
\end{proposition}

\begin{proof}
记 \(S_m:=\sum_{k=1}^m F_{k+1}=F_{m+3}-2\)。
则
\[
N(C(\omega))=\sum_{k=1}^m(1-\omega_k)F_{k+1}=S_m-N(\omega).
\]
又 \(F_{m+3}=F_{m+2}+F_{m+1}\) 给出
\(
S_m=F_{m+3}-2\equiv F_{m+1}-2\pmod{F_{m+2}}
\)，
从而
\[
r(C(\omega))\equiv (F_{m+1}-2)-r(\omega)=\tau_m(r(\omega))\pmod{F_{m+2}}.
\]
由定理 \ref{thm:fold-congruence-universal-property} 中的表述可取 \(\Fold_m(\omega)=w(r(\omega))\)，故
\[
\Fold_m(C(\omega))=w(r(C(\omega)))=w(\tau_m(r(\omega)))=\jmath_m(w(r(\omega)))=\jmath_m(\Fold_m(\omega)),
\]
得到 (1)。
对 (2)，由于 \(C\) 为 \(\Omega_m\) 上双射，
\[
d_m(\jmath_m(x))
=\card{\{\omega:\Fold_m(\omega)=\jmath_m(x)\}}
=\card{\{\omega:\Fold_m(C(\omega))=\jmath_m(x)\}}
=\card{\{\omega:\jmath_m(\Fold_m(\omega))=\jmath_m(x)\}}
=d_m(x),
\]
其中最后一步用到 \(\jmath_m\) 为对合。证毕。
\end{proof}

\begin{theorem}[平移 \texorpdfstring{$+F_{m+1}$}{+F_{m+1}} 诱导稳定类型的全长循环]\label{thm:fold-curvature-translation-full-cycle}
固定 \(m\ge 1\)，令 \(G:=\ZZ/(F_{m+2}\ZZ)\)，并沿用定理 \ref{thm:fold-congruence-universal-property} 的规范截面
\(
w:G\to X_m
\)。
在 \(G\) 上定义平移
\[
T_m:G\to G,\qquad T_m(r):=r+F_{m+1}\ (\mathrm{mod}\ F_{m+2}).
\]
则 \(T_m\) 为单一长循环：对任意 \(r\in G\)，其轨道长度恰为 \(|G|=F_{m+2}\)。
从而在稳定层上得到一条规范的全长回路
\[
w(r)\ \to\ w(T_m r)\ \to\ w(T_m^2 r)\ \to\ \cdots\ \to\ w(T_m^{F_{m+2}-1} r)\ \to\ w(r).
\]
\end{theorem}
\begin{proof}
由 \(\gcd(F_{m+1},F_{m+2})=1\)（Fibonacci 递推下的 Euclid 归纳）知 \(F_{m+1}\) 在循环群 \(G\) 中生成全体元素，故 \(r\mapsto r+F_{m+1}\) 的轨道长度为 \(|G|=F_{m+2}\)。
将该循环经双射 \(w\) 推前即得稳定层上的回路。证毕。
\end{proof}

\begin{theorem}[补码对称与曲率平移生成二面体作用]\label{thm:fold-curvature-dihedral-action}
固定 \(m\ge 1\)，沿用命题 \ref{prop:fold-complement-duality} 的对合
\(
\jmath_m:X_m\to X_m
\)
与余数对合 \(\tau_m(r)\equiv (F_{m+1}-2)-r\ (\mathrm{mod}\ F_{m+2})\)。
再令 \(\widetilde T_m:X_m\to X_m\) 为定理 \ref{thm:fold-curvature-translation-full-cycle} 的平移在稳定层上的搬运：
\[
\widetilde T_m(w(r)):=w(T_m r)\qquad(r\in G).
\]
则 \(\langle \jmath_m,\widetilde T_m\rangle\) 在 \(X_m\) 上给出一个忠实的二面体群作用，并满足关系
\[
\boxed{\;
\langle \jmath_m,\widetilde T_m\rangle\ \cong\ D_{F_{m+2}},
\qquad
\jmath_m\widetilde T_m\jmath_m=\widetilde T_m^{-1}. \;}
\]
\end{theorem}
\begin{proof}
由 \(\tau_m(r)=a-r\)（\(a:=F_{m+1}-2\)）可得 \(\tau_m^2=\mathrm{id}\)，故 \(\jmath_m^2=\mathrm{id}\)。
并且对任意 \(r\in G\)，有
\[
\tau_m(T_m r)\equiv a-(r+F_{m+1})\equiv (a-r)-F_{m+1}\equiv T_m^{-1}(\tau_m(r))\pmod{F_{m+2}},
\]
从而
\(
\jmath_m\widetilde T_m\jmath_m=\widetilde T_m^{-1}
\)。
另一方面，由定理 \ref{thm:fold-curvature-translation-full-cycle}，\(\widetilde T_m\) 的阶为 \(F_{m+2}\)。
因此生成关系与二面体群表示一致。
忠实性来自 \(w\) 为双射：若群元在 \(X_m\) 上作用平凡，则其在 \(G\) 上对 \(r\) 的作用亦平凡，从而群元为单位元。证毕。
\end{proof}

\endinput


